\section{Limitaciones y conclusiones}
\label{sec:conclusiones}

\textbf{Limitaciones.}
\begin{itemize}
  \item \textbf{Validación pequeña.} Con 58 dorados, los IC de la AUC macro miden unos
        $\pm0{,}06$, y los de cada hallazgo, más del doble. Por eso la mayoría de las
        ablaciones no es concluyente.
  \item \textbf{Techo del texto.} El techo lo pone la supervisión: el extractor tiene un 19\,\%
        de ruido, y en sinovitis y derrame el texto no recoge lo que se ve en la imagen.
  \item \textbf{Una sola vista y resolución gruesa.} Usamos una sola vista y mapas de
        7$\times$7, insuficientes para estructuras finas como los meniscos.
  \item \textbf{Umbral trasladado.} El umbral elegido en datos débiles se traslada mal a los
        dorados, que están enriquecidos.
  \item \textbf{Particiones.} El conjunto no trae identificadores de paciente, así que las
        particiones son por estudio.
  \item \textbf{Asignación de idioma.} El idioma, detectado automáticamente, es un proxy
        imperfecto del centro.
\end{itemize}

\textbf{Conclusiones.} Es posible entrenar un clasificador de doce hallazgos en RM de rodilla
sin etiquetas estructuradas, a partir de reportes en nueve idiomas.
\begin{itemize}
  \item La transferencia desde ImageNet supera a una CNN desde cero (0,715 frente a 0,654),
        siempre que el ajuste fino use una tasa de aprendizaje baja. Con tasa alta, el modelo
        colapsa a un predictor constante.
  \item Frente al ruido de las etiquetas, conservar toda la supervisión con parada temprana
        es mejor que descartar las etiquetas dudosas.
  \item Grad-CAM y el análisis por idioma muestran riesgos que la AUC macro oculta: atajos
        por correlación en el LCA y los meniscos, y una calidad de etiquetado desigual entre
        centros.
\end{itemize}
Como trabajo futuro proponemos fusionar las vistas sagital y coronal, recortar centrado en la
articulación para ganar resolución, sustituir las reglas por embeddings multilingües o un
modelo de lenguaje, y validar por centro antes de cualquier despliegue.
